\section*{\fontsize{16pt}{15pt}\selectfont DESIGN AND DEPLOYMENT OF METRO ETHERNET NETWORK USING MPLS \\ABSTRACT}

\begin{itemize}
    \item[-] \textbf{Research issues:}
    
    This project focuses on simulating a Metro Ethernet MAN infrastructure using Multiprotocol Label Switching (MPLS) technology to connect multiple enterprise branches with different internal LAN architectures.
    
    The report includes 5 chapters:
        \begin{itemize}
        \item[•] Chapter 1 - Overview and Theoretical Basis
        \item[•] Chapter 2 - System Design and Planning
        \item[•] Chapter 3 - Technical Configuration and Deployment
        \item[•] Chapter 4 - Experiment and Performance Evaluation
        \item[•] Chapter 5 - Result Analysis and Conclusion
        \end{itemize}
    \item[-] \textbf{Approach:}
        \begin{itemize}
        \item[•] \textit{Theory:} Researching Metro Ethernet standards, MPLS principles (Label Switching), OSPF and LDP routing protocols, and modern LAN architectures (Flat, 3-Tier, Leaf-Spine).
        \item[•] \textit{Practice:} Using Mininet and FRRouting to deploy the MPLS Backbone and customer networks, with automated performance measurement using Python scripts.
    \end{itemize}
    \item[-] \textbf{How to solve problems:}
    Constructing a simulation including Core (P), Edge (PE), and Customer (CE) routers. Configuring OSPF for routing and LDP for label distribution. Implementing three distinct LAN architectures at branches for performance evaluation.
    \item[-] \textbf{Results achieved:}
    The system successfully converged and performed accurate label switching. Automated tests provided data on throughput, delay, and jitter. Results showed that MPLS enhances packet processing speed and ensures high stability for inter-branch connections.
\end{itemize}

\newpage